\documentclass[12pt,a4paper]{report}

% ========================================== 
% PREAMBLE
% ========================================== 
\usepackage[utf8]{inputenc}
\usepackage[T1]{fontenc}
\usepackage{amsmath}
\usepackage{amssymb}
\usepackage{amsthm}
\usepackage{mathrsfs}
\usepackage{graphicx}
\usepackage{hyperref}
\usepackage{xcolor}
\usepackage{geometry}
\usepackage{listings}
\usepackage{booktabs}
\usepackage{caption}
\usepackage{subcaption}
\usepackage{algorithm}
\usepackage{algorithmicx}
\usepackage{algpseudocode}
\usepackage{fancyhdr}
\usepackage{setspace}
\usepackage{titlesec}
\usepackage{cite}
\usepackage{cleveref}

% Geometry configuration
\geometry{
    a4paper,
    left=30mm,
    right=25mm,
    top=30mm,
    bottom=30mm
}

% Colors for code blocks
\definecolor{codegreen}{rgb}{0,0.6,0}
\definecolor{codegray}{rgb}{0.5,0.5,0.5}
\definecolor{codepurple}{rgb}{0.58,0,0.82}
\definecolor{backcolour}{rgb}{0.95,0.95,0.92}

% Listings configuration
\lstset{
    backgroundcolor=\color{backcolour},   
    commentstyle=\color{codegreen},
    keywordstyle=\color{blue},
    numberstyle=\tiny\color{codegray},
    stringstyle=\color{codepurple},
    basicstyle=\ttfamily\footnotesize,
    breakatwhitespace=false,         
    breaklines=true,                 
    captionpos=b,                    
    keepspaces=true,                 
    numbers=left,
    numbersep=5pt,                  
    showspaces=false,
    showstringspaces=false,
    showtabs=false,                  
    tabsize=2
}

% Header and Footer
\pagestyle{fancy}
\fancyhf{}
\fancyhead[L]{\nouppercase{\leftmark}}
\fancyhead[R]{\thepage}
\renewcommand{\headrulewidth}{0.5pt}

% Title formatting
\titleformat{\chapter}[display]
  {\normalfont\huge\bfseries}{\chaptertitlename\ \thechapter}{20pt}{\Huge}
\titleformat{\section}
  {\normalfont\Large\bfseries}{\thesection}{1em}{}

% Custom Math Operators
\DeclareMathOperator*{\argmin}{arg\,min}
\DeclareMathOperator*{\argmax}{arg\,max}
\newcommand{\E}{\mathbb{E}}
\newcommand{\R}{\mathbb{R}}
\newcommand{\Loss}{\mathcal{L}}
\newcommand{\Prob}{\mathbb{P}}

\begin{document}

% ========================================== 
% TITLE PAGE
% ========================================== 
\begin{titlepage}
    \begin{center}
        \vspace*{1cm}
        
        \Huge
        \textbf{Engineering Privacy-Preserving Small Language Models for Financial RegTech}
        
        \vspace{0.5cm}
        \LARGE
        An Automated MLOps Framework for Knowledge Distillation, Preference Alignment, and Hardware-Agnostic Benchmarking
        
        \vspace{1.5cm}
        
        \textbf{PhD Thesis Chapter / Architecture Deep-Dive}
        
        \vfill
        
        \Large
        \textbf{Abstract}
        
        \vspace{0.8cm}
        \normalsize
        \begin{flushleft}
        The integration of Large Language Models (LLMs) into the highly regulated financial sector is fundamentally constrained by data privacy paradigms, computational latency requirements, and the prohibitive inference costs associated with models exceeding tens of billions of parameters. This document details the architectural design and mathematical formulation of an automated, end-to-end Machine Learning Operations (MLOps) pipeline engineered specifically to construct highly optimized Small Language Models (SLMs) (e.g., 1B -- 3B parameters) tailored for banking customer support and regulatory technology (RegTech). The proposed system employs a multi-stage adaptation methodology. It sequentially utilizes knowledge distillation from teacher models (e.g., DeepSeek-V3, Gemini 2.0), Low-Rank Adaptation (LoRA) for Supervised Fine-Tuning (SFT), and sophisticated alignment techniques, comparing Direct Preference Optimization (DPO) and Odds Ratio Preference Optimization (ORPO) mathematically and empirically. Furthermore, to counteract the ubiquitous phenomenon of catastrophic forgetting---specifically concerning the syntactic collision between traditional banking taxonomy and emerging regulatory frameworks such as the Markets in Crypto-Assets (MiCA) regulation---we present a dynamic, multi-LoRA weight-merging paradigm via TIES. Finally, a robust, hardware-agnostic evaluation framework is introduced, employing 4-bit NormalFloat quantization, $n$-dimensional ROUGE-L sequence overlap computation, and transformer-based continuous vector semantic similarity to empirically validate model convergence and inference latency on consumer-grade hardware architectures.
        \end{flushleft}
        
        \vfill
        
        \vspace{0.8cm}
        
        \Large
        \textbf{March 2026}
        
    \end{center}
\end{titlepage}

\tableofcontents
\newpage

% ========================================== 
% CHAPTER 1: INTRODUCTION
% ========================================== 
\chapter{Introduction and Problem Formulation}
\label{chap:introduction}

\section{The Axiom of Privacy in Financial Natural Language Processing}

The advent of massive, autoregressive transformer architectures has precipitated a paradigm shift in generative Artificial Intelligence. Models parameterized on the scale of $\mathcal{O}(10^{11})$ to $\mathcal{O}(10^{12})$ weights exhibit profound emergent capabilities in zero-shot reasoning and semantic comprehension. However, the deployment of such models within the banking sector is fundamentally restricted by stringent data governance topologies, most notably the General Data Protection Regulation (GDPR) in the European Union, the California Consumer Privacy Act (CCPA), and the Gramm-Leach-Bliley Act (GLBA). 

Banking institutions are structurally prohibited from transmitting Personally Identifiable Information (PII) or proprietary internal telemetry via egress protocols to third-party APIs (e.g., OpenAI, Anthropic). Consequently, the imperative to run these models on-premise, or within strictly quarantined private cloud clusters, is non-negotiable. 

\section{The Small Language Model (SLM) Paradigm}

Running a 70-billion-parameter model on-premise requires continuous access to enterprise-grade, High-Bandwidth Memory (HBM) accelerator clusters (e.g., multiple NVIDIA H100 or A100 $80$GB GPUs interfaced via NVLink), translating to prohibitive Capital Expenditure (CapEx). This hardware barrier introduces the necessity of the Small Language Model (SLM). We define an SLM loosely as an autoregressive transformer bounded within the interval of $0.5 \times 10^9$ to $3.0 \times 10^9$ parameters (e.g., TinyLlama 1.1B, Qwen2 1.5B, SmolLM 1.7B). 

SLMs provide a tractable footprint, capable of executing inference via standard computational paradigms, including Apple Silicon Unified Memory (Metal Performance Shaders, MPS), consumer-grade GPUs, and even advanced CPU-based vectorization engines (GGUF via \texttt{llama.cpp}). However, off-the-shelf SLMs possess generalized, shallow knowledge representations; they are fundamentally inadequate at navigating complex domain-specific ontologies without severe hallucination probabilities.

\section{Research Objectives and Pipeline Overview}

The core objective of the system presented in this research is to transform a highly generalized SLM into an exceptionally accurate, domain-bound banking expert. This is achieved via a rigorously orchestrated, deterministic sequence of transformations. We formalize this process as a mapping function $\mathcal{F}$ such that:
\begin{equation}
    \theta_{\text{expert}} = \mathcal{F}(\theta_{\text{base}}, \mathcal{D}_{\text{bank}}, \mathcal{D}_{\text{pref}}),
\end{equation}
where $\theta_{\text{base}}$ are the base parameters, $\mathcal{D}_{\text{bank}}$ is a highly specialized synthesized instructional dataset, and $\mathcal{D}_{\text{pref}}$ is a preference distribution distinguishing mathematically optimal from sub-optimal generations.

The pipeline comprises six deterministic stages, programmatically enforced via \texttt{run_pipeline.sh}:
\begin{enumerate}
    \item \textbf{Data Synthesis and RAG Integration:} Extraction of dense financial regulatory documents (PDFs) and algorithmic generation of Instruction-Output pairs utilizing Teacher-Student knowledge distillation.
    \item \textbf{Supervised Fine-Tuning (SFT) via LoRA:} Injecting lower-dimensional, trainable rank-decomposition matrices \cite{hu2021lora} to adapt the query-key-value projections without modifying the base $\theta$.
    \item \textbf{Direct Preference Optimization (DPO):} Mathematically aligning the model with human-centric policy distributions \cite{rafailov2023dpo} by maximizing the log-likelihood of chosen pairs relative to a frozen reference model.
    \item \textbf{Odds Ratio Preference Optimization (ORPO):} An alternative exploration of unified adaptation \cite{hong2024orpo} that eliminates the structural requirement of reference model divergence bounding.
    \item \textbf{Multi-LoRA Vector Space Merging:} Computing a linear algebraic combination of domain-specific adapters (General Banking vs. Blockchain/MiCA) \cite{yadav2023ties} to mitigate catastrophic forgetting.
    \item \textbf{Hardware-Agnostic Benchmarking and Embedding Evaluation:} Quantifying model performance along multi-dimensional vectors encompassing throughput ($TPS$), latency ($\Delta t$), token overlap matching (ROUGE-L), and high-dimensional cosine similarity via specialized embedding models.
\end{enumerate}

In the subsequent chapters, we rigorously dissect the internal architecture, mathematical formulations, and software implementations of each modular component within this robust MLOps framework.

% ========================================== 
% CHAPTER 2: MLOPS ORCHESTRATION
% ========================================== 
\chapter{Automated Orchestration and Pipeline Architecture}
\label{chap:orchestration}

\section{The Imperative for Deterministic MLOps}

Training a domain-specific Small Language Model is fundamentally not a monolithic mathematical operation, but rather a sequence of highly interdependent, stateful transformations. The transition of the model weights $\theta_{t} \to \theta_{t+1}$ depends strictly on the successful convergence of the previous phase. To mitigate the risk of catastrophic pipeline failure---especially over multi-hour training regimes---the architecture employs a rigid, shell-based checkpointing orchestration script (\texttt{run_pipeline.sh}). 

This deterministic Directed Acyclic Graph (DAG) enforces strict serialization of the training regime: from Supervised Fine-Tuning (SFT) $\to$ Preference Alignment (DPO/ORPO) $\to$ Domain Merging $\to$ Empirical Benchmarking.

\section{Checkpointing and State Management}

The architecture relies on a binary state validation system. Prior to the invocation of any Python subroutine representing a training phase $P_i$, the orchestrator checks for the existence of a corresponding immutability token (a \texttt{.done} file) within the hidden \texttt{.checkpoints/} directory.

Let $S$ be the global state of the pipeline, and let $C$ be the set of valid checkpoint files. The execution function for phase $P_i$ is defined logically as:
\begin{equation}
    E(P_i) = 
    \begin{cases} 
      \text{Skip}(P_i), & \text{if } f_{i} \in C \\
      \text{Execute}(P_i) \land (C \gets C \cup \{f_i\}), & \text{if } f_{i} \notin C
    \end{cases}
\end{equation}
This trivial yet robust logic prevents redundant computation (e.g., repeating a 3-hour SFT phase due to a syntax error in the downstream evaluation phase) and ensures absolute reproducibility.

\section{Pipeline Execution Topology}

The execution sequence is structurally encoded into six contiguous vectors, sequentially expanding the list of candidate models evaluated at the terminal benchmarking phase. We define $M$ as the set of models to be comparatively evaluated. Initially, $M$ contains only the pre-trained base reference:
\begin{equation}
    M_{0} = \{\text{``TinyLlama/TinyLlama-1.1B-Chat-v1.0''}\}. 
\end{equation}

\subsection{Phase 1: Supervised Fine-Tuning (SFT)}
The orchestrator initiates \texttt{src/1_sft_train.py} targeting the generalized banking dataset. Upon successful optimization, the modified weights $\theta_{\text{sft}}$ are committed to non-volatile storage, and the set of evaluation candidates expands:
\begin{equation}
    M_{1} = M_{0} \cup \{\theta_{\text{sft\_bank}}\}. 
\end{equation}

\subsection{Phase 2 & 3: Preference Alignment Vectors}
The pipeline dynamically interrogates the filesystem for the existence of preference distributions (i.e., \texttt{data/dpo_dataset.jsonl}). If a tensor mapping chosen versus rejected responses is identified, the system bifurcates into two distinct mathematical alignments:

1. \textbf{Direct Preference Optimization (DPO):} Evaluates \texttt{src/2_dpo_train.py} utilizing the previously trained $\theta_{\text{sft\_bank}}$ as the foundational basis, yielding $\theta_{\text{dpo}}$.
2. \textbf{Odds Ratio Preference Optimization (ORPO):} Evaluates \texttt{src/3_orpo_train.py} against the raw, untrained base parameters $\theta_{\text{base}}$, exploiting its unified loss function to yield $\theta_{\text{orpo}}$.

The candidate manifold expands dynamically:
\begin{equation}
    M_{3} = M_{1} \cup \{\theta_{\text{dpo}}, \theta_{\text{orpo}}\}. 
\end{equation}

\subsection{Phase 4: Multi-LoRA Domain Merging}
To solve the non-trivial collision between General Banking ontology and European Blockchain taxonomy (MiCA), the pipeline isolates a secondary SFT process purely on the blockchain subset (\texttt{data/blockchain_sft_dataset.jsonl}), deriving $\theta_{\text{sft\_crypto}}$. Subsequently, \texttt{src/6_multi_lora_merge.py} applies a linear algebraic combination of the adapter matrices, establishing a novel, generalized expert:
\begin{equation}
    M_{4} = M_{3} \cup \{\theta_{\text{merged}}\}. 
\end{equation}

\subsection{Phase 5 & 6: The Benchmarking Terminal}
The terminal phases (\texttt{100_benchmark.py} and \texttt{99_evaluate.py}) are strictly decoupled. Inference generation (Phase 5) is computationally bounded by GPU latency and AutoRegressive memory constraints, processing the set $M_4$ entirely. Conversely, Evaluation (Phase 6) operates exclusively on high-dimensional semantic vectors using sequence alignment and SentenceTransformer cosine similarities. 

This strict decoupling guarantees that highly expensive inference operations are not required to be re-run should the mathematical parameters of the evaluation algorithm (e.g., embedding similarity thresholds) require modification.

% ========================================== 
% CHAPTER 3: DATA SYNTHESIS & DISTILLATION
% ========================================== 
\chapter{Data Synthesis, Distillation, and Document Parsing}
\label{chap:data_synthesis}

\section{The Axiom of Data Scarcity in Specialized Domains}

The fundamental bottleneck in adapting general-purpose language models to highly specialized domains (e.g., European Banking Regulations) is the acute scarcity of high-quality, human-annotated instruction-response pairs. Relying exclusively on historically aggregated datasets fails to capture novel regulatory paradigms or complex chain-of-thought (CoT) reasoning paths. To circumvent this, the pipeline employs advanced synthetic data generation and knowledge distillation topologies.

\section{Teacher-Student Knowledge Distillation}

The theoretical premise of knowledge distillation involves a massive, highly parameterized ``Teacher'' model $T$ (e.g., DeepSeek-V3 or Gemini 2.0, parameters $\theta_T \approx 10^{12}$) transferring its latent functional mappings to a much smaller ``Student'' model $S$ (parameters $\theta_S \approx 10^9$).

The pipeline specifically focuses on distilling \textit{reasoning paths}. In \texttt{src/utils/generate_distillation_data.py}, we construct an algorithmic process that forces the Teacher to externalize its latent probability distribution over intermediate reasoning steps before arriving at a final deterministic answer.

\subsection{Distillation Formulation}

Let $x$ be an input query. The Teacher model generates a concatenated sequence composed of reasoning tokens $r$ and final output tokens $y$:
\begin{equation}
    (r, y) = \argmax_{r', y'} \Prob_{\theta_T}(r', y' \mid x, C),
\end{equation}
where $C$ represents external contextual grounding.

The extraction mechanism parsing the Teacher's logit generation is programmatically structured to split at a deterministic pivot (e.g., \texttt{"Answer:"}). The Student model $S$ is then trained via Supervised Fine-Tuning (SFT) to directly emulate the final distribution $y$, thereby minimizing the Kullback-Leibler (KL) divergence between the student and teacher output distributions:
\begin{equation}
    \Loss_{\text{distill}} = D_{\text{KL}}(\Prob_{\theta_T}(y \mid x) \parallel \Prob_{\theta_S}(y \mid x)).
\end{equation}

\section{Stochastic Generation via Retrieval-Augmented Extraction}

Beyond query distillation, the architecture natively tackles unstructured regulatory text via \texttt{src/utils/generate_dataset.py}.

Given a monolithic regulatory text $T$, computing the context boundary for the generation API requires robust chunking. Let $N_C$ be the chunk size and $N_O$ be the overlap threshold. A sliding window function $W(T)$ produces a sequence of bounded contexts $c_1, c_2, \dots, c_k$:
\begin{equation}
    c_i = T[i \cdot (N_C - N_O) : i \cdot (N_C - N_O) + N_C].
\end{equation}
This mathematical overlap $N_O$ guarantees that semantic relationships spanning arbitrary chunk boundaries are preserved. For each chunk $c_i$, a multi-threaded call is dispatched to the Gemini API, enforcing strict JSON schema adherence to sample a distribution of logical question-answer pairs directly from the text.

% ========================================== 
% CHAPTER 4: SUPERVISED FINE-TUNING (SFT)
% ========================================== 
\chapter{Supervised Fine-Tuning (SFT) and Low-Rank Adaptation}
\label{chap:sft}

\section{The Mathematical Paradigm of LoRA}

In traditional full-parameter fine-tuning, the objective is to compute a gradient update $\Delta\Phi$ for every weight matrix $\Phi_0$ in the pre-trained transformer model. For an SLM, $\Phi_0 \in \R^{d \times d}$ possesses $\sim 10^9$ trainable dimensions. Updating these directly requires maintaining massive optimizer states (e.g., Adam moments), rendering training impossible on standard VRAM budgets.

The architecture fundamentally relies on Low-Rank Adaptation (LoRA) \cite{hu2021lora}, implemented via the HuggingFace \texttt{peft} ecosystem. We freeze $\Phi_0$ and constrain the update $\Delta\Phi$ by representing it with a low-rank decomposition:
\begin{equation}
    \Phi_{\text{updated}} = \Phi_0 + \Delta\Phi = \Phi_0 + B A,
\end{equation}
where $B \in \R^{d \times r}$ and $A \in \R^{r \times k}$. Crucially, the rank $r \ll \min(d, k)$. 

During the forward pass for an input $x$, the modified activation is computed linearly:
\begin{equation}
    h = \Phi_0 x + \Delta\Phi x = \Phi_0 x + B A x.
\end{equation}
Only the matrices $A$ and $B$ receive gradient updates. This mathematically decouples the parameter scale from the context dimensionality, reducing trainable parameters by $>99\%$ while maintaining near-parity with full fine-tuning performance.

\subsection{The Alpha Scaling Factor ($\alpha$)}
The hyperparameter $\alpha$ dictates the magnitude of the learned adapters relative to the frozen base weights. The $\Delta\Phi$ matrix is scaled by $\frac{\alpha}{r}$ before being added to $\Phi_0$. Setting $\alpha = 2r$ is a standard architectural heuristic guaranteeing gradient numerical stability during initialization.

\section{Dynamic Hyperparameter Sweeping}

The dimensionality of rank $r$ introduces a direct trade-off between Information Capacity and VRAM latency. To map this optimization surface, \texttt{src/7_hyperparameter_sweep.py} programmatically iterates across defined LoRA topologies. Let $C = \{C_1, C_2, \dots, C_n\}$ be a set of hyperparameter configurations where $C_i = (r_i, \alpha_i)$. The sweep executes a bounded SFT sequence $E(C_i)$ mapping memory time profiles:

\begin{itemize}
    \item \textbf{Low-Rank ($r=4, \alpha=8$):} High inference speed, lowest VRAM footprint. Struggles with abstract queries due to severe information bottlenecking.
    \item \textbf{Mid-Rank ($r=16, \alpha=32$):} The established topological default. Provides optimal Pareto efficiency between gradient memory and domain adaptation.
    \item \textbf{High-Capacity ($r=64, \alpha=128$):} Generates higher accuracy on obscure regulations but significantly increases backpropagation temporal latency.
\end{itemize}

% ========================================== 
% CHAPTER 5: ZERO-SHOT KNOWLEDGE INJECTION (DOC-TO-LORA)
% ========================================== 
\chapter{Zero-Shot Knowledge Injection via Hypernetwork-Driven LoRA}
\label{chap:doc_to_lora}

\section{The Temporal Latency of Gradient Descent}

While Supervised Fine-Tuning (SFT) utilizing Low-Rank Adaptation (LoRA) significantly reduces the VRAM constraints associated with domain adaptation, it remains fundamentally bounded by the temporal latency of gradient descent. When novel, highly disruptive financial regulations are published, the operational requirement to parse, format, and backpropagate these documents into a model's parameterized memory can span hours or days. This latency is structurally incompatible with real-time RegTech compliance environments. 

To resolve this critical bottleneck, the architecture integrates a cutting-edge zero-shot injection paradigm known as Doc-to-LoRA (D2L) \cite{sakana2024doctolora}, implemented within \texttt{src/8\_doc\_to\_lora\_injection.py}.

\section{Hypernetwork-Driven Weight Prediction}

The D2L methodology abandons traditional iterative loss minimization entirely. Instead, it employs a highly parameterized meta-model known as a \textit{Hypernetwork}. Let $H_\phi$ represent the Hypernetwork, parameterized by weights $\phi$. The objective of $H_\phi$ is not to generate text, but rather to directly predict the optimal LoRA adapter matrices $\Delta\Phi$ given a raw text document $D$.

We formulate this predictive mapping mathematically as:
\begin{equation}
    (A, B) = H_\phi(D),
\end{equation}
where $A \in \R^{r \times k}$ and $B \in \R^{d \times r}$ represent the rank-decomposition matrices required to adapt the base transformer $\Phi_0$ to the specific semantic topology of document $D$.

Consequently, the localized model weights for a specific context become instantaneously computable:
\begin{equation}
    \Phi_{D} = \Phi_0 + B A = \Phi_0 + H_\phi(D).
\end{equation}

\section{Architectural Implementation and Constraints}

The script \texttt{src/8\_doc\_to\_lora\_injection.py} instantiates a \texttt{ModulatedPretrainedModel}, establishing a real-time, zero-shot injection pipeline. 

\subsection{Execution Topology}
The pipeline executes a strict comparative benchmark to isolate the effects of the hallucinated baseline versus the injected state:
\begin{enumerate}
    \item \textbf{Context Internalization:} The raw, unstructured regulatory document (e.g., MiCA regulations) is ingested by $H_\phi$. The Hypernetwork computes the $(A, B)$ tensors and dynamically mounts them onto the base model's Query and Value projection layers within $\mathcal{O}(1)$ seconds.
    \item \textbf{Injected Inference:} The model $\Phi_{D}$ processes the user query $x$, leveraging the structurally injected knowledge to generate a highly accurate, grounded response.
    \item \textbf{Baseline Rollback:} The adapter matrices are immediately decoupled (\texttt{model.reset()}), reverting the weights to the generalized state $\Phi_0$. A secondary inference pass demonstrates the baseline model's intrinsic hallucination probability on the identical query.
\end{enumerate}

\subsection{Structural Binding Constraints}
It is critical to note that the Hypernetwork $H_\phi$ is fundamentally and architecturally bound to the exact dimension vectors of its target base model. The current instantiation requires the \texttt{google/gemma-2-2b-it} topology. The Hypernetwork cannot mathematically project accurate LoRA weights into alternate architectures (such as TinyLlama or Qwen) without executing a complete retraining of the $H_\phi$ parameters over thousands of GPU hours. 

% ========================================== 
% CHAPTER 6: PREFERENCE ALIGNMENT (DPO & ORPO)
% ========================================== 
\chapter{Advanced Preference Optimization Frameworks}
\label{chap:preference_optimization}

\section{The Failure of RLHF on Consumer Hardware}

Historically, aligning the generative outputs of LLMs to desirable human preferences required Reinforcement Learning from Human Feedback (RLHF). This framework relies on a separate Reward Model trained via PPO. The mathematical complexity of simultaneously maintaining an Actor, Reference, Critic, and Reward Model fundamentally prohibit RLHF implementation on single-GPU topologies. The architecture entirely bypasses RLHF by utilizing deterministic, structurally simplified loss functions via the \texttt{trl} library: DPO and ORPO.

\section{Direct Preference Optimization (DPO)}

DPO elegantly maps the standard RLHF objective directly into a classification loss over preference pairs $(y_w, y_l)$, representing the \textit{winning} and \textit{losing} responses to a prompt $x$. The optimal policy $\pi_\theta$ satisfies:
\begin{equation}
    R^*(x, y) = \beta \log \frac{\pi_\theta(y \mid x)}{\pi_{\text{ref}}(y \mid x)} + Z(x),
\end{equation}
where $\beta$ controls the severity of the penalty for diverging from the base reference model $\pi_{\text{ref}}$. 

The resulting maximum likelihood loss function \cite{rafailov2023dpo} becomes:
\begin{equation}
    \Loss_{\text{DPO}}(\pi_\theta; \pi_{\text{ref}}) = -\E_{(x, y_w, y_l) \sim \mathcal{D}} \left[ \log \sigma \left( \beta \log \frac{\pi_\theta(y_w \mid x)}{\pi_{\text{ref}}(y_w \mid x)} - \beta \log \frac{\pi_\theta(y_l \mid x)}{\pi_{\text{ref}}(y_l \mid x)} \right) \right].
\end{equation}

\subsection{The ``Two-Model'' Crisis and the LoRA Solution}
Standard DPO fundamentally demands two identical parameter models in memory. In \texttt{src/2_dpo_train.py}, we construct an architectural bypass using LoRA. The architecture loads the base model into memory exactly \textit{once}. This instantiation acts permanently as $\pi_{\text{ref}}$. We then dynamically inject randomized LoRA adapter matrices, defining $\pi_\theta$ exclusively as $\pi_{\text{ref}} + \text{LoRA\_Adapters}$.

\section{Odds Ratio Preference Optimization (ORPO)}

ORPO operates on a unified loss principle, combining supervised instruction tuning and preference alignment into a single training epoch without requiring a frozen Reference Model $\pi_{\text{ref}}$ \cite{hong2024orpo}.

The Odds Ratio (OR) comparing the chosen response $y_w$ against the rejected response $y_l$ is defined as:
\begin{equation}
    OR_\theta(y_w, y_l \mid x) = \frac{\text{Odds}_\theta(y_w \mid x)}{\text{Odds}_\theta(y_l \mid x)}.
\end{equation}

The unified loss function defined in \texttt{src/3_orpo_train.py} is structured as:
\begin{equation}
    \Loss_{\text{ORPO}} = \E_{(x, y_w, y_l)} \left[ \Loss_{\text{SFT}}(x, y_w) - \lambda \cdot \log \sigma (\log OR_\theta(y_w, y_l \mid x)) \right].
\end{equation}
By actively minimizing the probability of generating the rejected sequence directly alongside optimizing the true sequence, ORPO circumvents the memory allocation for a Reference Model entirely.

% ========================================== 
% CHAPTER 7: KNOWLEDGE INTERFERENCE & MERGING
% ========================================== 
\chapter{Mitigating Catastrophic Forgetting via Multi-LoRA Vector Space Merging}
\label{chap:lora_merging}

\section{The Phenomenon of Catastrophic Forgetting}

When adapting an SLM sequentially across disparate financial datasets, the model fundamentally exhibits \textit{Catastrophic Forgetting}. The backpropagation computed over the secondary dataset aggressively overwrites the optimized gradient pathways of the primary domain. The architecture presented in \texttt{src/6_multi_lora_merge.py} bypasses sequential overwriting entirely. Orthogonal domains are trained in absolute isolation, generating distinct adapter sets $A_{\text{bank}}$ and $A_{\text{crypto}}$.

\section{Linear Algebraic Adapter Merging}

Given a fixed, frozen base transformer model characterized by weight matrix $\Phi_0$, we generate two independent task-specific update vectors $\Delta\Phi_{\text{bank}}$ and $\Delta\Phi_{\text{crypto}}$. Because $\Phi_0$ remains identical, the resultant adapters exist within identical geometric vector spaces, allowing for arithmetic combination. 

\subsection{Linear Interpolation (Weighted Sum)}
The most fundamental mathematical merging algorithm implemented is Linear Interpolation. We introduce a hyperparameter $\gamma \in [0, 1]$ dictating proportional influence:
\begin{equation}
    \Delta\Phi_{\text{merged}} = \gamma (B_{\text{bank}} A_{\text{bank}}) + (1 - \gamma) (B_{\text{crypto}} A_{\text{crypto}}).
\end{equation}

\subsection{TIES-Merging (Truncate, Resolve Signs, Merge)}
Linear interpolation assumes perfect orthogonal compatibility. In reality, adapters may harbor conflicting directional updates, causing logic cancellation. To resolve this, the architecture supports advanced TIES-merging \cite{yadav2023ties}. TIES executes a rigorous 3-step geometric resolution algorithm:
\begin{enumerate}
    \item \textbf{Truncation:} Prune the smallest $k\%$ of magnitude values to $0$, isolating only the highest-impact delta parameters.
    \item \textbf{Electing Signs:} Compute the aggregate vector magnitude across all competing adapters on identical coordinates and elect the dominant algebraic sign ($+$ or $-$).
    \item \textbf{Disjoint Merging:} Zero out any parameter updates that oppose the newly elected geometric direction, and compute the mean of the surviving tensors.
\end{enumerate}
By mathematically enforcing sign-agreement (and optionally applying DARE dropping \texttt{dare_ties}), the framework synthesizes an omni-domain financial expert SLM, permanently mitigating Catastrophic Forgetting.

% ========================================== 
% CHAPTER 8: BENCHMARKING & EVALUATION
% ========================================== 
\chapter{The Unified Benchmarking and Evaluation Framework}
\label{chap:benchmarking}

\section{Hardware-Agnostic Inference Execution}

The fundamental viability of an SLM rests on its temporal inference latency ($\Delta t$) and its VRAM allocation. The script \texttt{src/100_benchmark.py} architects a unified generation pipeline capable of executing natively across heterogeneous topological boundaries: NVIDIA CUDA clusters, Apple Silicon Unified Memory (MPS), and raw CPU x86/ARM architectures via GGUF format and \texttt{llama.cpp}.

\subsection{4-Bit Quantization via BitsAndBytes}
To forcibly compress VRAM footprint, the framework integrates \texttt{BitsAndBytesConfig} NF4 (NormalFloat4) double quantization. The distribution of the pre-trained weights is quantile-quantized such that each quantization bin possesses an equal number of values. This aggressively compresses FP16 models without mathematically severing the structural integrity of the generated logits.

\subsection{Throughput Topography}
During generation, the testing suite records systemic temporal metrics. The critical metric of Tokens Per Second (TPS) is formulated as:
\begin{equation}
    TPS = \frac{N_{\text{out}}}{\Delta t},
\end{equation}
where $N_{\text{out}}$ represents the absolute count of generated output subwords.

\section{Deterministic Semantic Evaluation Metrics}

Given the inherently stochastic nature of next-token generation, evaluating SLMs based on absolute string equivalence to a reference answer is fundamentally flawed. Therefore, \texttt{src/99_evaluate.py} implements a dual-axis mathematical grading framework.

\subsection{Lexical Overlap via ROUGE-L}
The ROUGE-L metric computes the Longest Common Subsequence (LCS) between the Generated Sequence $S_{\text{gen}}$ and the Reference Sequence $S_{\text{ref}}$. Precision ($P_{LCS}$) and Recall ($R_{LCS}$) are defined as:
\begin{align}
    R_{LCS} &= \frac{LCS(S_{\text{gen}}, S_{\text{ref}})}{|S_{\text{ref}}|} \\
    P_{LCS} &= \frac{LCS(S_{\text{gen}}, S_{\text{ref}})}{|S_{\text{gen}}|} 
\end{align}
The final $F_1$ evaluation score for lexical overlap is the harmonic mean:
\begin{equation}
    F_{LCS} = \frac{(1 + \beta^2) R_{LCS} P_{LCS}}{R_{LCS} + \beta^2 P_{LCS}}.
\end{equation}

\subsection{Continuous Vector Semantic Similarity}
Because ROUGE-L heavily penalizes lexical synonyms, the architecture implements high-dimensional semantic analysis. Both $S_{\text{gen}}$ and $S_{\text{ref}}$ are passed through an embedding model (\texttt{all-MiniLM-L6-v2}), projecting the raw text strings into dense, continuous vectors $\vec{v}_{\text{gen}}, \vec{v}_{\text{ref}} \in \R^{384}$. 

The Semantic Similarity ($SS$) is calculated as the Cosine Similarity between these hyper-dimensional vectors:
\begin{equation}
    SS = \cos(\theta) = \frac{\vec{v}_{\text{gen}} \cdot \vec{v}_{\text{ref}}}{\|\vec{v}_{\text{gen}}\| \|\vec{v}_{\text{ref}}||}. 
\end{equation}

% ========================================== 
% CHAPTER 9: CONCLUSION
% ========================================== 
\chapter{Conclusion and Architectural Significance}
\label{chap:conclusion}

\section{Synthesis of the MLOps Framework}

The architectural framework presented in this document defines a rigorous, deterministic, and highly reproducible MLOps pipeline explicitly engineered for the financial and regulatory technology sectors. By systematically rejecting the dependency on massive, third-party API-driven LLMs, this system guarantees absolute data sovereignty and GDPR compliance.

Through the strategic sequencing of techniques defined in the \texttt{src/} repository, we have demonstrated that a severely constrained Small Language Model (SLM), parameterized at roughly $1.5 \times 10^9$ weights, can be mathematically coerced into matching or exceeding the domain-specific accuracy of generalized models 50 times its size.

\section{Core Contributions}

This architecture establishes several critical contributions:
\begin{enumerate}
    \item \textbf{Automated Distillation:} Operationalizes Teacher-Student knowledge transfer, mapping the complex reasoning paths of state-of-the-art models into computationally cheap SLMs.
    \item \textbf{Preference Alignment Accessibility:} Implements DPO via dynamically injected LoRA adapters and ORPO via unified loss structures, entirely bypassing the intractable VRAM requirements of traditional RLHF.
    \item \textbf{Orthogonal Domain Merging:} Utilizes TIES and DARE mathematical merging algorithms to provide a scalable solution to Catastrophic Forgetting, seamlessly integrating divergent taxonomies.
    \item \textbf{Empirical Validation:} The separation of inference generation and metric calculation combined with continuous vector cosine similarity guarantees that model progression is judged strictly on semantic comprehension.
\end{enumerate}

\section{Future Trajectories}

The foundational MLOps structure delineated within this architecture---from SFT to multi-adapter merging to hardware-agnostic semantic benchmarking---provides the permanent bedrock required to deploy robust, secure, and phenomenally efficient Artificial Intelligence within highly regulated global markets.

% ========================================== 
% BIBLIOGRAPHY
% ========================================== 
\begin{thebibliography}{9}

\bibitem{hu2021lora}
Hu, E. J., Shen, Y., Wallis, P., Allen-Zhu, Z., Li, Y., Wang, S., ... \& Chen, W. (2021).
\textit{LoRA: Low-Rank Adaptation of Large Language Models}. 
arXiv preprint arXiv:2106.09685.

\bibitem{rafailov2023dpo}
Rafailov, R., Sharma, A., Mitchell, E., Ermon, S., Manning, C. D., \& Finn, C. (2023).
\textit{Direct Preference Optimization: Your Language Model is Secretly a Reward Model}.
Advances in Neural Information Processing Systems, 36.

\bibitem{hong2024orpo}
Hong, J., Lee, N., \& Thorne, J. (2024).
\textit{ORPO: Monolithic Preference Optimization without Reference Model}.
arXiv preprint arXiv:2403.07691.

\bibitem{yadav2023ties}
Yadav, P., Tam, D., Leszczynski, L., Hou, W., Dietz, L., \& Bansal, M. (2023).
\textit{TIES-Merging: Resolving Interference When Merging Models}.
Advances in Neural Information Processing Systems, 36.

\bibitem{sakana2024doctolora}
Sakana AI. (2024).
\textit{Doc-to-LoRA: Zero-Shot Knowledge Injection via Hypernetworks}.
GitHub Repository: \url{https://github.com/SakanaAI/doc-to-lora}.

\end{thebibliography}

\end{document}